\section{Optimised codebooks}
\label{app:codebooks}
\subsection*{ArMIS -- blind}
\begin{quote}\small
---

# Annotation Uncertainty Codebook for Arabic Tweets (Revised)



**Objective**

Judge the stability of the label (misogynistic or sexisitic) assignment for a given Arabic tweet. Estimate the likelihood that a human annotator would flip labels upon repeated review. Note: Linguistic features (dialect, typos) should not override semantic clarity.



**Dimension 1: Semantic Clarity of Hostility**

- **Indicator: Direct Attack** (Text explicitly attacks women/character, calling them corrupt/sinners/weak).

- **Indicator: Underlying Meaning** (Text refers to women but describes male perspective; e.g., "men don't like these women").

- **Rule:** Direct accusations of character or social failure by women decrease uncertainty. Ambiguity on whether the text attacks *women* or just *a situation involving women* increases uncertainty. *Downgrade linguistic noise (dialect/typos) weight; prioritize intent.*



**Dimension 2: Target Specificity**

- **Indicator: Generic Groups** (e.g., "sisters", "women", "all daughters").

- **Indicator: Specific Targets** (e.g., "this player", "this community").

- **Rule:** Generic groups increase uncertainty *only if* the attack relies on interpretation. If the text attacks "women in general" with clear misogyny (e.g., "women are corrupt"), uncertainty is **not uncertain**. Specific targets lower uncertainty.



**Dimension 3: Contextual Ambiguity**

- **Indicator: Satire/Irony** (Text mocking gender roles but unclear if the speaker supports or opposes the role).

- **Indicator: Single Focus** (Pure gender issues).

- **Indicator: Mixed Agenda** (Politics + gender).

- **Rule:** Satire increases uncertainty. Mixed agendas increase uncertainty *only if* the gender component is diluted or secondary to a non-gender topic. Pure focus decreases uncertainty.



**Dimension 4: Stylistic Noise (Secondary Factor)**

- **Indicator: Broken Syntax/Dialect** (Colloquialisms, typos).

- **Rule:** This raises uncertainty *only if* it obscures Dimension 1 (hostility). If the hostility is clear despite slang, uncertainty remains **not uncertain**.



**Combinatorial Logic**

- **not uncertain:** Clear hostility, specific or generic target (depending on clarity), standard or slang dialect (doesn't matter if intent is clear), single focus.

- **somewhat uncertain:** Intent is slightly masked (e.g., potential sarcasm), or generic target requires more context to confirm hostility, or minor linguistic ambiguity.

- **highly uncertain:** Intent is genuinely ambiguous (could be anti-men or neutral comment), heavy political context obscures gender focus, or text refers to women but it is unclear if they are the target of the action.



**Final Assessment**

not uncertain

---
\end{quote}
\subsection*{ArMIS -- labelled}
\begin{quote}\small
---

# Annotation Codebook: Uncertainty of Misogyny/Sexism Labels (Revised)



## Objective

The goal is to assess the **stability of the label assignment** by predicting how likely an annotator is to disagree with the assigned label when viewing the text again under different conditions. An instability rating reflects real-world annotator variance.



## Construct Dimensions and Indicators



### 1. Harm Specificity (Inherent vs. Situational)

This dimension measures whether the negative stance targets a woman's **inherent value** or **specific behaviors/choices**.



*   **Indicators of "Inherent Value" (Stable):**

    *   Text rhetoric that attacks women's intellect, soul, custody of opinion, or spiritual worth (e.g., "incomplete mind", "defective life").

    *   Text using common derogatory metaphors involving animals or objects to describe women (e.g., "cows", "bulls", "objects").

    *   Direct statements removing agency (e.g., "Allah takes women").

*   **Indicators of "Situational Behavior" (Unstable):**

    *   Text focusing on clothing, fashion, or social adaptation.

    *   Text discussing women in political or social contexts (e.g., "women in leadership", "women crying", "women voting").

    *   Text discussing religious physiology (e.g., menstruation) used as justification for role exclusion (high variance in acceptance).



### 2. Cultural Consensus (Dogma vs. Argument)

This dimension measures whether the text relies on a **widely accepted cultural trope** for abjection or presents a **specific argument**.



*   **Indicators of "Cultural Consensus" (Stable):**

    *   Phrasing that relies on common patriarchal religious interpretations (despite being pious-sounding) where "women < men" is the unstated norm.

    *   Standardized insults that rely on a shared cultural understanding of devaluation.

*   **Indicators of "Specific Argument" (Unstable):**

    *   Text attempting to reframe harm as "science", "politics", or "general observation".

    *   Text where the gendered harm is secondary to an event (e.g., a car crash, a political meeting).



### 3. Ambiguity of Intent (Attack vs. Complaint)

This dimension measures the separation between *criticism of a policy/behavior* and *criticism of the gender*.



*   **Indicators of "Attack" (Stable):**

    *   The subject is the gender (e.g., "women are...").

*   **Indicators of "Complaint" (Unstable):**

    *   The subject is an event where women are participants (e.g., "In politics... women usually lose").

    *   Text where the insult is indirect or requires heavy parsing to distinguish from neutral reporting.



## Decision Logic

Combine the dimensions below to assign the uncertainty level.



1.  **Assign "not uncertain" if:**

    *   The text has **Inherent Value Attack** AND/OR **Cultural Consensus**.

    *   The text explicitly uses standard insults, or religious/cultural tropes that universally map to devaluing women's worth in the local context, even if phrased religiously or metaphorically.

    *   *Reasoning:* Annotators will rarely disagree on standard insults or deeply entrenched hierarchies regarding gender worth.



2.  **Assign "somewhat uncertain" if:**

    *   The text has **Situational Behavior** OR **Specific Argument**.

    *   The text discusses gendered impacts on behavior, social roles, or specific events where harm is implied but not direct.

    *   *Reasoning:* Annotators may disengage and classify it as "neutral" or "political" rather than "harmful to gender".



3.  **Assign "highly uncertain" if:**

    *   The text has **Ambiguity of Intent** AND/OR **Low Cultural Consensus**.

    *   The text requires external context (idioms, heavy dialect, obscure slang) to determine if a woman is the target or the narrator.

    *   *Reasoning:* If the target or intent relies on decoding unknown context, annotators will lack stability.



**Override Rule:**

*   Do not label "not uncertain" if the text is a neutral description of an event (e.g., car crash) even if women are harmed (e.g. Example 6 is stable only because the *joke* was the attack, not the crash).

*   Do not label "somewhat uncertain" if the text is a direct attack on intellect/identity (e.g., "defective minds") even if it uses religious terminology.



---







not uncertain

---
\end{quote}
\subsection*{HSBrexit -- blind}
\begin{quote}\small
---

### Annotation Codebook: Judging Annotator Instability on Sensitive Issues



**Task**:

You are tasked with assessing the likelihood of disagreement among human annotators on a specific English tweet's classification (Hate Speech vs. Not Hate Speech). Focus strictly on the **stability of the label** across repeated annotator draws, not the objective correctness of the label itself. Identify if the text creates a scenario where reasonable annotators could assign different labels (Hate/Not Hate) based on standard hate speech guidelines.



**Scale Definitions**:

- **"not uncertain"**: The text's classification is the applicability of guidelines is clear. Political criticism, personal grievances, and specific individual targets are generally consistent across annotators. Consensus should be high.

- **"somewhat uncertain"**: The text contains political-contextual group generalizations or reports on sensitive events where the distinction between policy preference and identity attack is ambiguous. Consensus may vary.

- **"highly uncertain"**: The text obscures the author/source or relies heavily on coded language/idioms unique to political subcultures where annotation guidelines do not provide a clear distinction criterion.



**Dimension 1: Nature of Negative Sentiment**

- **Indicator**: Does the negative sentiment target a group's inherent identity/stereotypes or a specific policy preference?

- **Low Uncertainty (Not Uncertain)**: Language criticizes *actions*, *policies*, or *outcomes* (e.g., "want foreigners out," "policy is bad").

- **High Uncertainty (Somewhat/Highly)**: Language assigns negative *attributes* or *characterization* to a protected group (e.g., "Stupid Brits," "proud Muslims"). If the group is characterized by a deficit (stupid, evil, dangerous) rather than an action (wants policy X), uncertainty increases.



**Dimension 2: Target Scope and Tone**

- **Indicator**: Who is the recipient of the sentiment?

- **Low Uncertainty (Not Uncertain)**: Specific grievance (e.g., "I was searched"), Direct Address ("You are an X [derogatory]"), or Clear Political Opinion without generalization (e.g., "The government failed").

- **High Uncertainty (Somewhat/Highly)**: Broad targeting of a protected identity using political rhetoric (e.g., "Brits," "Immigrants," "Muslims" used as generalizations for policy).



**Dimension 3: Contextual Transparency**

- **Indicator**: Is the source clear? Is it reporting on the event or the event itself?

- **Low Uncertainty (Not Uncertain)**: Original post, clear ownership of the opinion.

- **High Uncertainty (Somewhat/Highly)**: News digest reporting on hate acts (e.g., "Journalist called X because of Y"), as annotators may disagree on if the *report* is the hate or the *context* implies it.



**Decision Rules**:

1.  **Payment vs. Attack (Identity)**: If the text makes a policy claim ("Foreigners should leave") -> **"not uncertain"**. If the text makes an identity claim ("Brits are stupid" / "Foreigners are a threat") -> **"somewhat uncertain"**.

2.  **Direct Address**: If the text addresses a specific target ("You are an immigrant") as a direct insult -> **"not uncertain"**. Personal grievances are also **"not uncertain"**.

3.  **News Reporting**: If the tweet reports on an incident involving hate speech or groups in a news format (e.g., "#BBC Journalist") -> **"somewhat uncertain"** (Ambiguity on if the broadcast is the hate).

4.  **No Auto-Upgrade on Repost**: Reposts are only marked **"highly uncertain"** if the *content* of the repost is itself ambiguous or the original intent is lost. Do not flag an original grievance simply because it is retweeted (**Ref:** Ex 4).

5.  **Default for Political Generalizations**: Do not default to "not uncertain" for political generalizations about groups (e.g., "Foreigners," "Nationalities"). Assign **"somewhat uncertain"** unless the generalization is strictly about a vote/policy choice (e.g., "I voted for X" -> Not Uncertain).



**Instruction for Execution**:

If the text relies on political sentiment attached to a group identity (rather than policy actions), mark **"somewhat uncertain"**. If the target is a specific policy preference, a personal grievance, or a direct insult, mark **"not uncertain"**. Only mark **"highly uncertain"** if the author/source is ambiguous or the language is coded. Do not use "highly uncertain" for standard political debates unless the line between speech and hate is non-existent (e.g. extreme slurs).

---
\end{quote}
\subsection*{HSBrexit -- labelled}
\begin{quote}\small
An English tweet about Brexit and immigration was labelled for hate speech.

The annotator could assign only these labels:

- not hate speech

- hate speech



Judge how uncertain that label is for this item.
\end{quote}